\subsubsection{Gateway (Go relay სერვერი)}

Gateway dumb relay ტიპის სერვერია: ის დოკუმენტის შესახებ
არაფერს ინახავს, არც სნეფშოტს, არც ოპერაციების ისტორიას არც CRDT
დოკუმენტს. მისი საქმეა ოპერაციების კლიენტებთან სწრაფად გადაგზავნა. 

მისი ძირითადი ცნებებია:
\begin{itemize}[itemsep=0pt]
  \item \textbf{Hub}, ოთახების რეესტრი: ქმნის ოთახებს (\texttt{POST
        /rooms}), პოულობს მათ იდენტიფიკატორით და შლის სესიის ბოლოს;
  \item \textbf{Room}, ერთი სესიის სივრცე: წევრების სია და გადაგზავნის
        ლოგიკა. მიღებული ფრეიმი ეგზავნება ოთახის ყველა წევრს გამგზავნის
        გარდა;
  \item \textbf{Client}, Room-ის წევრი კლიენტის სტრუქტურა. 
\end{itemize}

მიუხედავად dumb როლისა, გეითვეი სერვერს რამდენიმე აქტიური მოვალეობა მაინც
აქვს. ახალი წევრის შემოსვლისას ის ჰოსტს სნეფშოტის მოთხოვნის საკონტროლო
ფრეიმს უგზავნის და ჰოსტის პასუხს მხოლოდ ახალ წევრს გადასცემს, ასე
გვიან შემოსული სტუმარი სწრაფად იღებს დოკუმენტის მიმდინარე მდგომარეობას,
ისე, რომ სერვერს დოკუმენტის შენახვა არ სჭირდება. წევრის ყოველ
შემოსვლა-გასვლაზე ოთახი ყველა Peer-ს ატყობინებს და ახალ წევრს, არსებული წევრებზე სრულ ინფორმაციას გადასცემს.
გეითვეი ასევე ინახავს Permission data-ს კლიენტებზე და კლიენტებისგან შემოსული data-ზე ახდენს შესაბამის ვალიდაციებს:
არ აკეთებს ისეთი ფრეიმების დაგზავნას, რომლის გამომგზავნსაც write პრივილეგია არ აქვს.

უსაფრთხოებისა და მდგრადობის მხრივ: ყველა HTTP ენდპოინტი (გარდა
\texttt{/health} და \texttt{/ws}) მოითხოვს bearer ტოკენს, რომელიც მხოლოდ
ჰოსტის აპლიკაციამ იცის; HTTP endpoint-ები დაცულია IP Rate Limiter-ით. 
